% Part:sets-functions-relations
% Chapter: sets
% Section: reduction

\documentclass[../../../include/open-logic-section]{subfiles}

\begin{document}

\olfileid{sfr}{siz}{red}

\olsection{归约}

\begin{editorial}
  本节通过归约证明不可数性，对应 \olref[nen]{sec} 中的结果。
  另一种稍简的版本对应 \olref[nen-alt]{sec} 中的结果，见 \olref[red-alt]{sec}。
\end{editorial}

我们之前通过对角线论证证明了 $\Pow{\PosInt}$ 是不可数的。
我们也已经证明了 $\Bin^\omega$（所有由 $0$ 和 $1$ 组成的无限序列的集合）是不可数的。
证明 $\Pow{\PosInt}$ 不可数还有另一种方法：
证明\emph{如果 $\Pow{\PosInt}$ 是可数的，那么 $\Bin^\omega$ 也是可数的}。
既然我们知道 $\Bin^\omega$ 不可数，那么 $\Pow{\PosInt}$ 也不可能可数。
这被称为将一个问题\emph{归约}（reducing）到另一个问题——在此例中，我们将枚举 $\Bin^\omega$ 的问题归约为枚举 $\Pow{\PosInt}$ 的问题。
后者的解（即 $\Pow{\PosInt}$ 的一个枚举）将给出前者的解（即 $\Bin^\omega$ 的一个枚举）。

如何将枚举集合~$B$ 的问题归约为枚举集合~$A$ 的问题？
我们需要提供一种将~$A$ 的枚举转化为~$B$ 的枚举的方法。
最简单的方法是定义一个满射函数 $f\colon A \to B$。
如果 $x_1$, $x_2$, \dots{} 枚举了~$A$，那么 $f(x_1)$, $f(x_2)$, \dots{} 就枚举了~$B$。
在我们的情形下，需要寻找一个满射函数 $f\colon \Pow{\PosInt} \to \Bin^\omega$。

\begin{prob}
证明：如果存在单射函数 $g\colon B \to A$，且 $B$~不可数，则 $A$~也不可数。
请通过说明如何利用~$g$ 将~$A$ 的枚举转化为~$B$ 的枚举来证明这一点。
\end{prob}

\begin{proof}[{\olref[nen]{thm:nonenum-pownat}} 的归约证明]
假设 $\Pow{\PosInt}$ 是可数的，因而存在它的一个枚举 $Z_{1}$, $Z_{2}$, $Z_{3}$, \dots。

定义函数 $f \colon \Pow{\PosInt} \to \Bin^\omega$ 如下：$f(Z)$ 为序列 $s_{k}$，其中 $s_{k}(n) = 1$ 当且仅当 $n \in Z$，否则 $s_k(n) = 0$。
这显然定义了一个函数，因为只要 $Z \subseteq \PosInt$，任意 $n \in \PosInt$ 要么属于 $Z$，要么不属于 $Z$。
例如，正偶数集 $2\PosInt = \{2,
4, 6, \dots\}$ 被映射到序列 $010101\dots$，空集被映射到 $0000\dots$，而集合 $\PosInt$ 本身被映射到 $1111\dots$。

该函数也是满射的：每个由 $0$ 和 $1$ 组成的序列都对应某个正整数集合，即由序列中取值为~$1$ 的位置所对应的整数组成的集合。
更确切地说，假设 $s \in \Bin^\omega$。定义 $Z \subseteq \PosInt$ 为：
\[
Z = \Setabs{n \in \PosInt}{s(n) = 1}
\]
那么 $f(Z) = s$，这可以根据~$f$ 的定义验证。

现在考虑列表
\[
f(Z_1), f(Z_2), f(Z_3), \dots
\]
因为 $f$ 是满射的，$\Bin^\omega$ 中的每个成员都必定作为~$f$ 在某个自变元处的值出现，因此也必定出现在该列表中。
所以这个列表必然枚举了所有的~$\Bin^\omega$。

因此，如果 $\Pow{\PosInt}$ 是可数的，那么 $\Bin^\omega$ 就是可数的。
但是 $\Bin^\omega$ 是不可数的 (\olref[nen]{thm:nonenum-bin-omega})。
从而 $\Pow{\PosInt}$ 是不可数的。
\end{proof}

\begin{explain}
人们很容易混淆归约的方向。
例如，满射函数 $g \colon \Bin^\omega \to B$ 并\emph{不能}由此证明 $B$ 是不可数的。
（考虑由 $g(s) = s(1)$ 定义的 $g \colon \Bin^\omega \to \Bin$，该函数将 $0$ 和 $1$ 的序列映射为其第一个元素。它是满射的，因为有些序列以 $0$ 开头，有些以 $1$ 开头。但 $\Bin$ 是有限的。）
还要注意，函数~$f$ 必须是满射的，否则论证无法成立：
那样就无法保证 $f(x_1)$, $f(x_2)$, \dots{} 包含~$B$ 的所有元素。
例如，\[
h(n) = \underbrace{000\dots0}_{\text{$n$ $0$的}}
\]
定义了一个函数 $h\colon \PosInt \to \Bin^\omega$，但 $\PosInt$ 却是可数的。
\end{explain}

\begin{prob}
通过归约论证，证明所有\emph{正整数对集合}的集合是不可数的。
\end{prob}

\begin{prob}\label{sfr:siz:red:prob:nat-nat}
  通过归约论证，证明所有函数 $f\colon \Nat \to \Nat$ 构成的集合~$X$ 是不可数的。
  （提示：给出一个从 $X$ 到~$\Bin^\omega$ 的满射函数。）
\end{prob}

\begin{prob}
通过归约论证，证明 $\Nat^\omega$（自然数的无限序列构成的集合）是不可数的。
\end{prob}

\begin{prob}
设 $P$ 为从正整数集到集合 $\{0\}$ 的所有函数的集合，设 $Q$ 为从正整数集到集合 $\{0\}$ 的所有\emph{偏}函数的集合。
证明 $P$~是可数的而 $Q$~不是。
（提示：将枚举 $\Bin^\omega$ 的问题归约为枚举~$Q$ 的问题）。
\end{prob}

\begin{prob}
设 $S$ 为所有从正整数集到集合 $\{0,1\}$ 的满射函数的集合，即 $S$ 由所有满射的~$f\colon \PosInt \to \Bin$ 组成。证明 $S$ 是不可数的。
\end{prob}

\begin{prob}
证明由所有实数构成的集合~$\Real$ 是不可数的。
\end{prob}

\end{document}